\documentclass[12pt]{article}
\usepackage[utf8]{inputenc}
\usepackage[T1]{fontenc}
\usepackage{amsmath,amssymb}
\usepackage{geometry}
\geometry{margin=2.5cm}

\begin{document}

\noindent ord$(a)=n$. \ $\Rightarrow \langle a\rangle = \{e,a,\ldots,a^{n-1}\}$.

\noindent $\Rightarrow$ ord$(a) = |\langle a\rangle|$

\noindent când $|G|<\infty$

\bigskip
\noindent\fbox{Teoremă}: \ Fie $G$ = un grup abelian finit.

\[
S=\{\text{ord}(a) \mid a\in G\} \subseteq \mathbb{N}^*.
\]

\begin{enumerate}
\item[(1)] $n\in S$ şi $d\mid n \Rightarrow d\in S$.
\item[(2)] $m,n\in S$ şi $(m,n)=1 \Rightarrow mn\in S$.
\item[(3)] $m,n\in S \Rightarrow [m,n]\in S$.
\item[(4)] $m=\max S \Rightarrow$
\begin{enumerate}
\item[(a)] $x^m=e$ $(\forall)\,x\in G$.
\item[(b)] $m,n$ au aceiaşi divizori, $n=$ord$\,G$
\item[(c)] $x^n=e$ $(\forall)\,x\in G \Rightarrow m\mid n$.
\end{enumerate}
\end{enumerate}

\bigskip
\noindent\underline{Dem}:

\noindent (1) şi (2) -- ca la Lema anterioară.\footnote{formularea exactă de aici e neclară în original -- se pare că trimite spre demonstraţia deja făcută a acestor două proprietăţi.}

\bigskip
\noindent (3). \ Fie $a,b\in G$ a.î.\ ord $a=m$, ord $b=n$.

\[
m=p_1^{\alpha_1}\cdots p_r^{\alpha_r}
\]
\[
n=p_1^{\beta_1}\cdots p_r^{\beta_r}
\]
\[
f_i=\max\{\alpha_i,\beta_i\}
\]

\noindent $\exists\, c_i\in G$ a.î.\ ord $c_i=p_i^{f_i}$ $(\forall)\,i=\overline{1,r}$.

\bigskip
\noindent $p_i^{f_i}\mid m$ sau $n$\footnote{acest pas, care justifică de ce $p_i^{f_i}\in S$ (fiind un divizor al lui $m$ sau al lui $n$, aplicând proprietatea (1)), e compactat în original.}

\[
\big(p_i^{f_i},\,p_j^{f_j}\big)=1 \ \Rightarrow \ \text{ord}(c_ic_j) = p_i^{f_i}p_j^{f_j}
\]

\bigskip
\noindent Prin inducţie -- ord$(c_1\cdots c_r) = p_1^{f_1}\cdots p_r^{f_r} = [m,n]$

\bigskip
\noindent (4) Fie $g\in G$ a.î.\ ord $g=m$

\begin{enumerate}
\item[(a)] ord$(a)\mid m$ $(\forall)\,a\in G$.
\end{enumerate}

\end{document}
